\section{如何写一个靠谱的 evaluation report}

如果评估不是只看一个数字，那报告也不应该只写一个数字。一个好的 evaluation report 至少应该说明：
\begin{itemize}
    \item 任务定义和输入来源。
    \item 调用方式和任何 scaffolding。
    \item 指标定义、聚合方式和样本数量。
    \item 置信区间、方差或重复运行结果。
    \item 明确的局限与已知失败模式。
\end{itemize}

\begin{figure}[H]
\centering
\begin{tikzpicture}[node distance=0.9cm, every node/.style={font=\small, align=center}]
\node[draw, rounded corners, fill=blue!10, minimum width=4.0cm, minimum height=0.8cm] (r1) {task definition};
\node[draw, rounded corners, fill=blue!18, below=of r1, minimum width=4.0cm, minimum height=0.8cm] (r2) {prompt / system setup};
\node[draw, rounded corners, fill=blue!26, below=of r2, minimum width=4.0cm, minimum height=0.8cm] (r3) {metric + aggregation};
\node[draw, rounded corners, fill=blue!34, below=of r3, minimum width=4.0cm, minimum height=0.8cm] (r4) {confidence + variance};
\node[draw, rounded corners, fill=blue!42, below=of r4, minimum width=4.0cm, minimum height=0.8cm] (r5) {limitations};
\draw[-{Latex[length=3mm]}, thick] (r1) -- (r2);
\draw[-{Latex[length=3mm]}, thick] (r2) -- (r3);
\draw[-{Latex[length=3mm]}, thick] (r3) -- (r4);
\draw[-{Latex[length=3mm]}, thick] (r4) -- (r5);
\end{tikzpicture}
\caption{一个靠谱的 evaluation report 应该显式写出实验条件，而不是只给单个分数}
\end{figure}

\begin{importantbox}{什么叫“报告写完整了”}
不是把所有数字都堆出来，而是让读者知道这个数字是在什么规则下得到的、它能否复现、它对你的目标是否真的有意义。
\end{importantbox}

\subsection{置信区间不是装饰}

如果样本量很小，单次分数的波动可能很大。尤其是 pairwise 或少样本任务，最好报告重复实验或置信区间，而不是把某一次运行当成定论。

\begin{knowledgebox}{如何读置信区间}
如果两个模型分数差距很小，而误差条大量重叠，那“领先”未必有实际意义。报告里如果完全不提方差，通常说明作者希望你忽略不确定性。
\end{knowledgebox}

\begin{warningbox}{最常见的报告错误}
只报平均值，不报样本数；只报总体分数，不报失败案例；只报正结果，不报对照组。这些都会让评估看起来比它实际更稳。
\end{warningbox}

\subsection{一个最小检查清单}

拿到一份新的 benchmark 结果时，你至少可以快速检查这几项：
\begin{enumerate}
    \item 它测的是知识、推理、对话、工具还是安全？
    \item 它是测模型、系统还是方法？
    \item 输入和 prompt 是否固定？
    \item 输出是否依赖 judge？
    \item 有没有重复运行和方差？
    \item 有没有已知污染或饱和问题？
\end{enumerate}

\begin{importantbox}{读 benchmark 的底线}
如果这些问题答不出来，就不要急着根据分数下结论。先把实验条件搞清楚，再看模型到底改进了什么。
\end{importantbox}

